\documentclass[11pt]{article}

\usepackage{amsmath,amssymb,amsfonts}
\usepackage{graphicx}
\usepackage{textcomp}
\usepackage[utf8]{inputenc}
\usepackage{tabularx}
\usepackage{array}
\usepackage{multirow}
\usepackage[T1]{fontenc}
\usepackage{booktabs} 
\usepackage{amssymb} 
\usepackage{geometry}
\geometry{a4paper, margin=1.5cm}
\begin{document}

\Large

\noindent{\bf Questionnaire \hspace{5cm} Code:}
\medskip\hrule \medskip 

\noindent{\bf Age: \hspace{7cm} Class:}
\newline
\newline
\noindent{\textbf{Gender:} $\bigcirc$ Female \quad $\bigcirc$ Male \quad $\bigcirc$ Other \quad $\bigcirc$ Prefer not to answer }
\newline

\noindent\textbf{Instructions:} Imagine being in the situations described below. Evaluate each situation in terms of how much anxiety you would feel during the specific event, putting an \textbf{X} in the column that corresponds to your level of anxiety. Remember that there are no right or wrong answers.
\newline
\renewcommand{\arraystretch}{1.5}
\setlength{\tabcolsep}{4pt} % horizontal spacing in cells

\noindent
\resizebox{\textwidth}{!}{%
\begin{tabular}{|p{7cm}|c|c|c|c|c|}
\hline
\multicolumn{1}{|c|}{\textbf{Situation}} & \multicolumn{5}{c|}{\textbf{Level of anxiety}} \\
\cline{2-6}
\multicolumn{1}{|c|}{} & \textbf{Very little} & \textbf{Little} & \textbf{Moderate} & \textbf{Quite a lot} & \textbf{A lot} \\
\hline
Using multiplication tables and schemes to do math exercises &  &  &  &  &  \\
\hline
Thinking about the math test scheduled for tomorrow &  &  &  &  &  \\
\hline
Paying attention to the teacher solving a difficult math operation on the board &  &  &  &  &  \\
\hline
Taking a written math test &  &  &  &  &  \\
\hline
Doing many difficult math exercises at home for the next lesson &  &  &  &  &  \\
\hline
Paying attention during math class &  &  &  &  &  \\
\hline
Watching a classmate solve a math exercise &  &  &  &  &  \\
\hline
Being questioned “by surprise” in math &  &  &  &  &  \\
\hline
Dealing with a new math topic &  &  &  &  &  \\
\hline
\end{tabular}}

\newpage

\noindent\textbf{Instructions:} Evaluate how good you think you are at handling the different situations described below. Put an \textbf{X} in the column that corresponds to your level of ability. Remember that there are no right or wrong answers. 
\newline

\noindent\begin{tabular}{|p{7cm}|c|c|c|c|c|}
  \hline
  \textbf{How good are you at...} & \textbf{Not at all} & \textbf{A little} & \textbf{Medium} & \textbf{Quite} & \textbf{Very} \\
  \hline
  Learning mathematics & & & & & \\
  \hline
  Solving math problems & & & & & \\
  \hline
  Mental calculation & & & & & \\
  \hline
  Solving operations in columns & & & & & \\
  \hline
  Multiplication tables & & & & & \\
  \hline
  Completing homework assigned for home & & & & & \\
  \hline 
  Focusing on studying math without being distracted & & & & & \\
  \hline
  Committing to math study when you have other interesting things to do & & & & & \\
  \hline
  Organizing school activities & & & & & \\
  \hline
  Being interested in school subjects & & & & & \\
  \hline
\end{tabular}

\newpage

\noindent{\bf Questionnaire \hspace{5cm} Code:}
\medskip\hrule \medskip 
\noindent\textbf{Instructions:} Evaluate how much you agree with the statements below by putting an \textbf{X} in the box that corresponds to your opinion. Remember that there are no right or wrong answers. 
\newline

\noindent\begin{tabular}{|p{7cm}|c|c|c|c|c|}
  \hline
  \textbf{How much do you agree?} & \textbf{Not at all} & \textbf{A little} & \textbf{Medium} & \textbf{Quite} & \textbf{Very} \\
  \hline
  I like studying mathematics & & & & & \\
  \hline
  Mathematics is useful in everyday life & & & & & \\
  \hline
  I would rather not have math lessons & & & & & \\
  \hline
  I find mathematics interesting & & & & & \\
  \hline
  I think mathematics will also be useful for my future & & & & & \\
  \hline
  Mathematics is a subject that motivates me to do my best & & & & & \\
  \hline 
  I don’t feel involved during math lessons & & & & & \\
  \hline
  I would like to go deeper into mathematics beyond what we do at school & & & & & \\
  \hline
\end{tabular}

\newpage

\noindent\textbf{Instructions:} Evaluate how much you agree with the statements below by putting an \textbf{X} in the box that corresponds to your opinion. Remember that there are no right or wrong answers. 
\newline

\noindent\begin{tabular}{|p{7cm}|c|c|c|c|c|}
  \hline
  \textbf{How much do you agree?} & \textbf{Not at all} & \textbf{A little} & \textbf{Medium} & \textbf{Quite} & \textbf{Very} \\
  \hline
  I worry that my classmates may think I have less math ability because of my gender & & & & & \\
  \hline
  When I get a bad grade in math I fear that it confirms the stereotype about my gender & & & & & \\
  \hline
  I think girls should not deal with mathematics because it is not suitable for them & & & & & \\
  \hline
  Sometimes I think I shouldn’t be good at math because of my gender & & & & & \\
  \hline
  I worry that my math teacher may think I’m not good because of my gender & & & & & \\
  \hline
  In mathematics, boys have a natural talent that girls cannot make up for with effort & & & & & \\
  \hline 
  To achieve the same results in mathematics, boys need to put in less effort & & & & & \\
  \hline
  Girls have to work harder than boys to achieve the same results in mathematics & & & & & \\
  \hline
  Despite effort, girls remain on average less good at mathematics than boys & & & & & \\
  \hline
\end{tabular}

\end{document}
